\section{Octree} \label{sec:octree}
The spatial domain can be subdivided into a hierarchical tree structure in many ways. For 3D problems, the most common approach is to use an Octree, which offers many opportunities to optimize algorithmic efficiency. An Octree is a tree data structure where each node is bisected along the Cartesian $x$, $y$, and $z$ axes, producing 8 equally sized child nodes, and these child nodes are then further subdivided into 8 nodes until a certain threshold is reached.

\subsection{Nodes of the Octree} \label{subsec:node}

We define an \texttt{FMMNode} data structure to represent each node in the Octree, which contains the following information:
\begin{description}
    \item [\textbf{Node geometry}:] The struct keeps track of the node's bounding box, its center, and its radius. It also stores the integer grid coordinates that represent the node's position in the Octree.
    \item [\textbf{Tree Hierarchy}:] Each node is assigned a node index, which serves as the node's unique identifier within a global list of all the nodes. The node also stores its tree level and the index of its parent node. Most importantly, it stores the indices of its 8 child nodes, which are used to navigate the tree.
    \item [\textbf{Mesh and Interaction data}:] The node stores an iterator for the particles contained within it. It maintains a neighbor list that contains the indices of neighboring nodes. These interactions must be computed directly. It also stores a list of far-field nodes, which are sufficiently far away to be approximated by multipole expansions.
    \item [\textbf{Mathematical expansion data}:] The node stores \texttt{P}m which is the order used for the mutlitpole expansion, and the number of coefficients required based on the order. It also stores the multipole and local expansion coefficients, which are used to approximate the potential and forces from distant nodes.
    \item [\textbf{FMM Translation Operators}:] The node stores the necessary translation operators for the FMM algorithm. 
    \begin{enumerate}
        \item \ac{P2M} Matrix: Translates raw particles (mesh sources) to a node's Multipole expansion.
        \item \ac{M2M} Matrix: Translates a child's Multipole expansion up to its parent's Multipole expansion.
        \item \ac{M2L} Matrices: Converts a distant node's Multipole expansion into this node's Local expansion.
        \item \ac{L2L} Matrix: Translates a parent's Local expansion down to its child's Local expansion.
        \item \ac{L2P} Matrix: Evaluates the Local expansion directly onto the target particles (mesh cells).
    \end{enumerate}
\end{description}

\todo{add a figure of Octree}

\subsection{Splitting the domain to make an Octree} \label{subsec:octree_splitting}
The underlying spatial data-structue sets the stage for the computational efficiency of the FMM algorithm. It dictates how the source-target interactions are categorized as either near-field, needing direct computation, or far-field, and thus suitable for approximation via multipole expansions. An adaptive, heirarchical Octree allows for a systematic way to manage these interactions.

We start by defining the root node as cube the encompasses the entire geometry. 
We  iterate over the vertices of all the mesh elements to find the minimum and maximum coordinates.
We then add a 1\% padding to the bounding box as a safety margin to ensure that all elements are contained within the root node, and to avoid edge cases where elements lie exactly on the boundary of the node. 
This safety margin can help guard against floating point precision issues occurring precisely at the boundary, which could lead to elements on or near the boundary being misclassified as outside the node.
This root node is then subdivided into 8 child nodes, and the elements are assigned to the child nodes based on their centroids. This is repeated recursively for each child node until the number of elements left in a node are less than a predefined threshold, at which point the node is designated as a leaf node. This threshold is a critical parameter that can be manually set at runtime. 


% The Octree is constructed by a top-down recursive subdivision of the spatial domain. 
\subsubsection{Space-Filling Curves and Morton Encoding} \label{subsubsec:morton-encoding}

We map the geometric centers of the boundary elements from three-dimensional space into a one-dimensional array using a space-filling curve, specifically, the Morton Z-curve. Morton encoding converts the 3D coordinates of each mesh element into a single integer by interleaving the bits of the binary representations of the x, y, and z coordinates, generating that element's Morton key. This ordering assigns numerically adjacent Morton keys to spatially adjacent elements. This allows us to sort the elements into a linearly ordered array,  where adjacent elements get clustered together, improving cache locality. This cache locality helps reduce the overhead of assigning the elements to the Octree's nodes. While the Morton Z-order curve may not be the optimum space-filling curve for 3D problems, its simplicity of implementation via bitwise operations lends itself to efficient execution \cite{malhotrapvfmmparallelkernel2015}.


\section{Building the Translation Matrices} \label{sec:translation-matrices}

The core optimization of our FMM implementation is the precomputation of the translation matrices. These matrices are used to perform the various translations required by the FMM algorithm, such as \ac{P2M}, \ac{M2M}, \ac{M2L}, \ac{L2L}, and \ac{L2P}. We compute the matrices once upfront and store them in the \texttt{FMMNode} data structure for each node in the Octree. This allows us to avoid expensive redundant calculations during the FMM algorithm, significantly improving the computational efficiency of our implementation.

We group the discussion of the translation matrices by the stage of the tree traversal they are used in. This traversal process involves mathematically translating the s